\documentclass[12pt]{article}
\usepackage[utf8]{inputenc}
\usepackage[T1]{fontenc}
\usepackage[margin=1in]{geometry}
\usepackage{setspace}
\usepackage{xcolor}
\usepackage[colorlinks=true, urlcolor=blue!60!black, linkcolor=blue!60!black]{hyperref}
\setlength{\parindent}{0em}
\setlength{\parskip}{0.6em}
\renewcommand{\baselinestretch}{1.0}
\begin{document}
\begin{center}
{\LARGE \textit{Troxel v. Granville}, 530 U.S. 57 (2000)\par}
\vspace{0.5em}
{\large Case Brief --- Two-Page Edition\par}
\end{center}

\section{Memory Jogger}
A court applying a ``breathtakingly broad'' visitation statute to override a fit mother's limits on grandparent visits violated her fundamental parental right (\href{https://www.law.cornell.edu/supremecourt/text/99-138}{Cornell LII}).

\section{Detailed Case Facts}
After their son Brad died, Jenifer and Gary Troxel sought more visitation with his daughters than the girls' mother, Tommie Granville, would allow; she offered one day a month. A Washington statute let ``any person'' petition at any time for visitation serving the child's best interest. Though no one questioned Granville's fitness, the trial court ordered a weekend a month and a summer week, presuming visitation was best.

\section{Procedural History}
The Washington Supreme Court held the statute facially unconstitutional because it required no showing of harm and swept too broadly.

\section{Judicial Votes}
Affirmed, 6--3, with no majority opinion. O'Connor wrote for Rehnquist, Ginsburg, and Breyer; Souter and Thomas concurred in the judgment; Stevens, Scalia, and Kennedy dissented separately.

\section{Holding}
Parents have a fundamental right to direct the care, custody, and control of their children. A court must give a fit parent's decision special weight; overriding it under a bare best-interest standard, with the presumption reversed, was unconstitutional as applied. Whether harm must be shown was left open.

\section{Analysis of Opinions}
\textbf{Plurality:} Granville was fit, the trial court presumed against her, ignored her offer of visitation, and rested on thin findings---a mere disagreement over what was better. \textbf{Souter:} the statute was invalid on its face. \textbf{Thomas:} strict scrutiny applies, and the State lacked even a legitimate interest. \textbf{Stevens:} children have their own interests in established bonds. \textbf{Scalia:} judges should not enforce unenumerated rights against state law. \textbf{Kennedy:} the harm rule was wrong, and de facto parents differ from strangers.

\section{Examples: Future Applications}
\textbf{Same-Side:} a law letting any neighbor or teacher obtain overnight visits whenever a judge finds them enriching is invalid. \textbf{Opposite-Side:} visitation ordered after a parent is adjudicated unfit is valid. \textbf{Fence-Sitter:} a former partner who co-raised a child for seven years seeks visitation without alleging harm.

\section{Critique}
Emily Buss called parental rights ``adrift in the middle,'' since the plurality defined neither the standard of review nor special weight. Child-welfare scholars warn that rigid parental autonomy slights children's relationships in blended and kinship families.

\section{Key Quotations}
\begin{quote}
``The liberty interest at issue in this case---the interest of parents in the care, custody, and control of their children---is perhaps the oldest of the fundamental liberty interests recognized by this Court.''
\end{quote}
\noindent --- Justice O'Connor, plurality opinion.

\end{document}
